\usepackage{sidecap}
\usepackage{quiver}

\usepackage{xcolor}
\definecolor{burgundy}{rgb}{0.5, 0.0, 0.13}

\usepackage{verbatim}
\usepackage{lmodern}
\usepackage{amsmath}
\usepackage{amssymb} 
\usepackage{amsthm} 
\usepackage{thmtools}
\usepackage{enumitem}
\usepackage{graphicx}
\usepackage[normalem]{ulem}
\usepackage[hidelinks]{hyperref}%hidelinks avoid the red boxes around the hyperlinks
\hypersetup{
  colorlinks = true,
  linktoc = all,
  linkcolor = burgundy, 
  urlcolor = Gray,
  citecolor = burgundy,
  pdfencoding = unicode,
  hypertexnames = false,
}

% --- keep long references inside the margins ---
\usepackage{xurl}                   % let URLs break at any character
\setlength{\emergencystretch}{3em}  % last-resort stretch for otherwise unbreakable lines



\usepackage[ 
style           	= alphabetic, 
sorting 			= nyt,
url					= false,
citetracker     	= true, 
natbib          	= true, 
hyperref      		= true,
backend      		= biber,
maxbibnames         = 99,
maxcitenames        = 2,
giveninits          = true,
uniquename    		= init,
parentracker  	  	= true,
backref        	 	= true,
backrefstyle    	=three]{biblatex}
\AtBeginBibliography{\Urlmuskip=0mu\relax} % no stretched gaps inside broken URLs
\setlength{\bibitemsep}{0.4\baselineskip} % Sets spacing between references
% Uses an underline when an author has multiple publications (in line with the AMS style)
\renewcommand*{\bibnamedash}{\underline{\hspace{3em}}\kern 0.1em} 
% Uses a semicolon delimeter for multiple citations
\renewcommand{\multicitedelim}{\addsemicolon\space}
% Removes the "p." when just a number is referenced
\DeclareFieldFormat{postnote}{#1}
% Uniformly italicizes title of all works
\DeclareFieldFormat[article,inbook,incollection,inproceedings,patent,thesis,unpublished]{title}{\textit{#1\isdot}} 
\DeclareFieldFormat{journaltitle}{#1}  % Makes Journal title in normal text
% Takes away the "ser." in front of Series Title that is part of the ieee styling
\DeclareFieldFormat[book,inbook,incollection,inproceedings]{series}{#1} 

\renewbibmacro*{in:}{\setunit{\addcomma\space}}


\usepackage{mathtools}
\usepackage{indentfirst}
\usepackage{tabularx}
\usepackage{bbm}
\usepackage[capitalise]{cleveref}
\usepackage{stmaryrd}
\usepackage{tabularx}
\usepackage{mathrsfs}
%\usepackage{mathabx}
%\usepackage{showframe}

\usepackage[left=1.1in,top=2cm,right=1.1in,bottom=2cm]{geometry}
%\usepackage[left=1.2in,top=2cm,right=1.2in,bottom=2cm]{geometry}

\makeatletter
\newcommand{\mylabel}[2]{#2\def\@currentlabel{#2}\label{#1}}
\makeatother

\setlength{\marginparwidth}{2cm}

\usepackage{array}
\newcolumntype{C}[1]{>{\centering\let\newline\\\arraybackslash\hspace{0pt}}m{#1}}


\usepackage{tikz} %pictures
\usepackage{tikz-cd} %diagrams
\newcommand{\btk}{\begin{tikzcd}}
\newcommand{\etk}{\end{tikzcd}}
\usetikzlibrary{arrows,calc,decorations.markings,bending}
\usepackage{xparse}



\usepackage{blindtext}


\makeatletter
\newcommand{\@bbify}[1]{
  \ifcsname b#1\endcsname
  \message{WARNING: Overwriting b#1 with blackboard letter!}
  \fi
  \expandafter\edef\csname b#1\endcsname
  {\noexpand\ensuremath{\noexpand\mathbb #1}\noexpand\xspace}}
\newcommand{\@calify}[1]{
  \ifcsname c#1\endcsname
  \message{WARNING: Overwriting c#1 with calligraphic letter!}
  \fi 
  \expandafter\edef\csname c#1\endcsname
  {\noexpand\ensuremath{\noexpand\mathcal #1}\noexpand\xspace}}
\newcommand{\@bfify}[1]{
  \ifcsname bf#1\endcsname
  \message{WARNING: Overwriting c#1 with bold letter!}
  \fi
  \expandafter\edef\csname bf#1\endcsname
  {\noexpand\ensuremath{\noexpand\mathbf #1}\noexpand\xspace}}
% Now we define all the commands
\newcounter{@letter}\stepcounter{@letter}
\loop\@bbify{\Alph{@letter}}\@calify{\Alph{@letter}}\@bfify{\Alph{@letter}}
\ifnum\the@letter<26\stepcounter{@letter}\repeat
\makeatother

\newenvironment{tz}{\begin{center}\begin{tikzpicture}}{\end{tikzpicture}\end{center}}
\tikzstyle{d}=[double distance=.3ex]
\tikzstyle{w}=[preaction={draw=white,-,line width=5pt}]

\newcounter{diagram}
\renewcommand{\thediagram}{\thetheorem}
\newenvironment{diagram}{
\setcounter{diagram}{\value{theorem}}
\refstepcounter{theorem}\refstepcounter{diagram}
\begin{center}
\normalfont{(\thediagram)}\hfill\begin{tikzpicture}[baseline=(current bounding box.center)]}
    {\end{tikzpicture}\hfill\text{ }\end{center}}

\tikzset{%
node distance=1.5cm, la/.style={scale=0.8}, lasmall/.style={scale=0.75}, over/.style={auto=false,fill=white,inner sep=1.5pt, minimum size=0, outer sep=0},
    symbol/.style={%
        draw=none,
        every to/.append style={%
            edge node={node [sloped, allow upside down, auto=false]{$#1$}}},
            
    }, pro/.style={postaction={decorate,decoration={
        markings,
        mark=at position .5 with {\node at (0,0) {$\bullet$};}
      }},
      inner sep=.9ex,
      },
      prosmall/.style={postaction={decorate,decoration={
        markings,
        mark=at position .5 with {\node at (0,0) {$\scriptstyle \bullet$};}
      }},
      inner sep=.9ex,
      },
  n/.style={double equal sign distance, -implies}, t/.style={double distance=2.5pt, -implies, postaction={draw,-}},
}

\newcommand{\pushout}[1]{\node at ($({#1})-(10pt,-10pt)$) {$\ulcorner$};}

\newcommand{\centered}[1]{\begin{tabular}{l} #1 \end{tabular}}



\newcommand{\Set}{\mathrm{Set}}
\newcommand{\colim}{\mathrm{colim}}


\newcommand{\Cof}{\mathrm{Cof}}
\newcommand{\an}{\ensuremath{\mathrm{An}}}
\newcommand{\trfib}{\ensuremath{\mathrm{TrFib}}}
\newcommand{\cof}{\ensuremath{\mathrm{cof}}}
\newcommand{\cellu}{\ensuremath{\mathrm{cell}}}
%\newcommand{\an}{\cof(\J)}
\newcommand{\nfib}{\ensuremath{\mathrm{NFib}}}
%\newcommand{\nfib}{\inj(\J)}
\newcommand{\fib}{\ensuremath{\mathrm{fib}}}
\newcommand{\proj}{\ensuremath{\mathrm{proj}}}


\newcommand{\inj}{\mathrm{inj}}
\newcommand{\Hom}{\mathrm{Hom}}


\newcommand{\fibrant}{\mathrm{fib}}
\newcommand{\Path}{\mathrm{Path}}
\newcommand{\id}{\mathrm{id}}
\newcommand{\ob}{\mathrm{ob}}
\newcommand{\op}{\mathrm{op}}
\newcommand{\Sq}{\mathrm{Sq}}
\newcommand{\cat}{\mathrm{Cat}}
\newcommand{\Wf}{\mathcal{W}_f}
\newcommand{\dblcat}{\mathrm{DblCat}}
\newcommand{\twocat}{2\mathrm{Cat}}
\newcommand{\ttwocat}{\mathbf{2}\cat}
\newcommand{\twogpd}{2\mathrm{Gpd}}
\newcommand{\bbH}{\mathbb{H}}



\renewcommand{\bfF}{\mathbf{F}}


\newcommand{\dblgrid}{\mathrm{DblGr}_{1-\mathrm{Id}}}
\newcommand{\dbldersch}{\mathrm{DblDerSch}}


\newcommand{\Eadj}{E_\mathrm{adj}}







\newlist{rome}{enumerate}{7}
\setlist[rome]{label=(\roman*)}


% Theorem-like environments
\newtheorem{thm}{Theorem}[section]
\newtheorem{cor}[thm]{Corollary}
\newtheorem{prop}[thm]{Proposition}
\newtheorem{lem}[thm]{Lemma}
\newtheorem{claim}[thm]{Claim}
\newtheorem*{conjecture}{Conjecture}

% Special theorem numbering
\declaretheorem[name=Theorem,numbered=yes]{theoremA}
\renewcommand{\thetheoremA}{\Alph{theoremA}}

\numberwithin{equation}{thm}
% Definition-style environments
\theoremstyle{definition}
\newtheorem{defn}[thm]{Definition}
\newtheorem{terminology}[thm]{Terminology}
\newtheorem{ex}[thm]{Example}
\newtheorem{notation}[thm]{Notation}
\newtheorem{descr}[thm]{Description}
\newtheorem{constr}[thm]{Construction}
\newtheorem{paragr}[thm]{}

% Remark-style environments
\theoremstyle{remark}
\newtheorem{rmk}[thm]{Remark}
\newtheorem{obs}[thm]{Observation}
%%%%%%%%%%


% Cleveref customization
\crefname{section}{Section}{Sections}
\crefname{thm}{Theorem}{Theorems}
\crefname{cor}{Corollary}{Corollaries}
\crefname{prop}{Proposition}{Propositions}
\crefname{lem}{Lemma}{Lemmata}
\crefname{claim}{Claim}{Claims}
\crefname{defn}{Definition}{Definitions}
\crefname{terminology}{Terminology}{Terminologies}
\crefname{ex}{Example}{Examples}
\crefname{notation}{Notation}{Notations}
\crefname{descr}{Description}{Descriptions}
\crefname{constr}{Construction}{Constructions}
\crefname{rmk}{Remark}{Remarks}
\crefname{obs}{Observation}{Observations}
\crefname{paragr}{Paragraph}{Paragraphs}

%========================================================================================
% TODO NOTES AND COMMENTS
%========================================================================================
\setlength{\marginparwidth}{2.5cm}
\usepackage[
  textwidth = 2.5cm,
  textsize = small, 
  colorinlistoftodos
]{todonotes}

% Comment bubble commands
\newcommand{\msnote}[1]{\todo[color=red!40,linecolor=red!40!black,size=\tiny]{#1}}
\newcommand{\msnoteil}[1]{\ \todo[inline,color=red!40,linecolor=red!40!black,size=\normalsize]{#1}}

\newcommand{\pvnote}[1]{\todo[color=orange!40,linecolor=orange!40!black,size=\tiny]{#1}}
\newcommand{\pvnoteil}[1]{\ \todo[inline,color=orange!40,linecolor=orange!40!black,size=\normalsize]{#1}}

\newcommand{\cbnote}[1]{\todo[color=blue!40,linecolor=blue!40!black,size=\tiny]{#1}}
\newcommand{\cbnoteil}[1]{\ \todo[inline,color=blue!40,linecolor=blue!40!black,size=\normalsize]{#1}}

\newcommand{\dtnote}[1]{\todo[color=green!40,linecolor=green!40!black,size=\tiny]{#1}}
\newcommand{\dtnoteil}[1]{\ \todo[color=green!40,linecolor=green!40!black,size=\normalsize]{#1}}

\newcommand{\stnote}[1]{\todo[color=cyan!40,linecolor=cyan!40!black,size=\tiny]{#1}}
\newcommand{\stnoteil}[1]{\ \todo[color=cyan!40,linecolor=cyan!40!black,size=\normalsize]{#1}}
\newcommand{\stnoteinline}[1]{\todo[color=cyan!40,linecolor=cyan!40!black,size=\normalsize, inline]{#1}}

\newcommand{\jnnote}[1]{\todo[color=yellow!40,linecolor=yellow!40!black,size=\tiny]{#1}}
\newcommand{\jnnoteil}[1]{\ \todo[inline,color=yellow!40,linecolor=yellow!40!black,size=\normalsize]{#1}}

%========================================================================================


% Cofibrations, fibrations, and weak equivalences
\DeclareMathOperator{\Fib}{Fib}
\newcommand{\cofM}{\Cof_\cM}
\newcommand{\fibM}{\Fib_\cM}
\newcommand{\weM}{\cW_\cM} % Assuming \cW is defined elsewhere or used directly
\newcommand{\trivcofM}{\cofM \cap \weM}
\newcommand{\trivfibM}{\fibM \cap \weM}

% Generating sets
\newcommand{\IM}{\cI_\cM}
\newcommand{\JM}{\cJ_\cM}
\newcommand{\IN}{L\cI}
\newcommand{\JN}{L\cJ}

% Cell complexes and injectivity
% NOTE: Use \text{-...} for hyphenated words in math mode.
\newcommand{\INcof}{\IN\text{-cof}}
\newcommand{\INinj}{\IN\text{-inj}}
\newcommand{\JNcof}{\JN\text{-cof}}
\newcommand{\JNinj}{\JN\text{-inj}}

% Notational variants for N
\newcommand{\cofN}{\cC}
\newcommand{\fibN}{\cF}
\newcommand{\weN}{\cW}
\newcommand{\trivcofN}{\cofN \cap \weN}
\newcommand{\trivfibN}{\fibN \cap \weN}


%========================================================================================
% FIBRATIONS
%========================================================================================


\DeclareMathOperator{\Nfib}{NFib}
\DeclareMathOperator{\DiscFib}{DiscFib}
\DeclareMathOperator{\splitfib}{split}
\DeclareMathOperator{\Ho}{Ho}

%========================================================================================
% CATEGORY TYPES 
%========================================================================================

\newcommand{\scat}{\mathrm{sCat}}
\newcommand{\bicat}{\mathrm{Bicat}}
\newcommand{\Bicat}{\mathrm{Bicat}}
\newcommand{\Top}{\mathrm{Top}}
\newcommand{\sSet}{\mathrm{sSet}}
\newcommand{\cSet}{\mathrm{cSet}}
\newcommand{\Psh}{\mathrm{Psh}}
\newcommand{\Cat}{\mathrm{Cat}}


% Slice categories
\newcommand{\slicetwocat}{\twocat_{/\cC}}
\newcommand{\Csharp}{\cC^{\sharp}}
\newcommand{\mslicetwocat}{\twocat^{+}_{/ \Csharp}}

%========================================================================================
% SIMPLICIAL SETS AND NERVES
%========================================================================================

% Basic simplicial sets
\newcommand{\Deltacat}{\mathrm{\Delta}}
\newcommand{\Dop}{\Deltacat^{\mathrm{op}}} % More semantic than hard-coding Delta


% Scaled and marked simplicial sets
\newcommand{\scaled}{\mathrm{sc}}
\newcommand{\markedscaled}{\mathrm{ms}}
\newcommand{\mbold}{\mathrm{mb}}
\newcommand{\sbold}{\mathrm{s}}

\newcommand{\scsSet}{\mathrm{sSet}^{\scaled}}
\newcommand{\mssSet}{\mathrm{sSet}^{\markedscaled}}
\newcommand{\mbssSet}{\mathrm{sSet}^{\mbold}}

%========================================================================================
% MARKED-SCALED STRUCTURES (No changes needed)
%========================================================================================

\newcommand{\markedscat}[1]{\left({#1}, E^1_{#1},E^2_{#1}\right)}
\newcommand{\markedscsset}[1]{\left({#1}, E_{#1},T_{#1}\right)}
\newcommand{\markedduskinnerve}[1]{\left(N_2({#1}), E_{#1},T_{#1}\right)}


%========================================================================================
% SPECIAL OBJECTS AND CONSTRUCTIONS
%========================================================================================
% NOTE: Use \DeclareMathOperator for function-like names to ensure correct spacing.
\DeclareMathOperator{\core}{Core}
\DeclareMathOperator{\Arr}{Ar}
\DeclareMathOperator{\target}{cod}

%========================================================================================
% MATH OPERATORS
%========================================================================================

% Hom functors
\DeclareMathOperator{\Aut}{Aut}
\DeclareMathOperator{\End}{End}
\DeclareMathOperator{\Fun}{Fun}
\DeclareMathOperator{\PsFun}{PsFun}
  \DeclareMathOperator{\Pathtwo}{\bO_{2}}
    \DeclareMathOperator{\Map}{Map}
\newcommand{\Funlex}{\Fun^{\lex}}
\newcommand{\HOM}{\ensuremath{\textup{\textsc{Hom}}}}

% (Co)limits
\newcommand{\cofib}{\ensuremath{\textup{cofib}}}
\DeclareMathOperator{\coker}{coker}
\DeclareMathOperator*{\hocolim}{hocolim}
\DeclareMathOperator*{\holim}{holim}

% General operators

\DeclareMathOperator{\interval}{\mathbb{I}}
\DeclareMathOperator{\Sing}{Sing}
\DeclareMathOperator{\Lan}{Lan}
\newcommand{\co}{\mathrm{co}}
\newcommand{\coop}{\mathrm{coop}}
\newcommand\partialbf{{\pmb{\partial}}}
\newcommand{\sk}{\operatorname{Sk}}
\newcommand{\An}{\operatorname{An}}
\newcommand{\im}{\operatorname{im}}
\renewcommand{\ker}{\operatorname{ker}}
\renewcommand{\Im}{\operatorname{Im}}
\renewcommand{\Re}{\operatorname{Re}}



\newcommand{\set}[1]{\left\lbrace #1 \right\rbrace }


% Custom label commands
\newcommand{\stlabel}[1]{{\upshape{(\ref*{#1}.\arabic*)}}}
\newcommand{\enumref}[2]{(\hyperref[#1.#2]{\ref*{#1}.\ref*{#1.#2}})}
\newcommand{\enumrefalt}[2]{\hyperref[#1.#2]{\ref*{#1}.\ref*{#1.#2}}}

\newcommand{\walkingadjoint}{E_{\mathrm{adj}}}
\newcommand{\isoname}[1]{\overset{#1}{\cong}}


%Adjunctions

\newcommand{\adjto}[4]{
	\begin{tikzcd}[ampersand replacement = \&]
		{#1:#2} \& {#3:#4}
		\arrow[""{name=0, anchor=center, inner sep=0}, shift left=1.3, from=1-1, to=1-2]
		\arrow[""{name=1, anchor=center, inner sep=0}, shift left=1.3, from=1-2, to=1-1]
		\arrow["\dashv"{anchor=center, rotate=-90}, draw=none, from=0, to=1]
	\end{tikzcd}}

\newcommand{\adjtorev}[4]{
	\begin{tikzcd}[ampersand replacement = \&]
		{#1:#2} \& {#3:#4}
		\arrow[""{name=0, anchor=center, inner sep=0}, shift left=1.3, from=1-1, to=1-2]
		\arrow[""{name=1, anchor=center, inner sep=0}, shift left=1.3, from=1-2, to=1-1]
		\arrow["\vdash"{anchor=center, rotate=-90}, draw=none, from=0, to=1]
	\end{tikzcd}}	

\newcommand{\adjtoRelo}[2]{
	\begin{tikzcd}[ampersand replacement = \&]
		{#1} \& {#2}
		\arrow[""{name=0, anchor=center, inner sep=0}, shift left=1.3, from=1-1, to=1-2]
		\arrow[""{name=1, anchor=center, inner sep=0}, shift left=1.3, from=1-2, to=1-1,hook']
		\arrow["\dashv"{anchor=center, rotate=-90}, draw=none, from=0, to=1]
	\end{tikzcd}}
  

    \newcommand{\kerodontag}[1]{\href{http://kerodon.net/tag/#1}{Tag #1}} 

%    \newcommand{\weN}{\cW}




%%%%%%%%%%%%%%%%%%%%%%%%%%
\NewDocumentCommand{\celli}{ O{} O{n} O{\cellslide} O{\celllength} m m m }{
  \coordinate (mid) at ($({#5})!{#3}!({#6})$);
  \coordinate (start) at ($(mid)!{#4}!({#5})$);
  \coordinate (end) at ($(mid)!{#4}!({#6})$);
  \draw[#2] (start) to node(label)[inner sep=4pt,outer sep=0,minimum size=0,#1]{{#7}} (end);
   \coordinate (far) at ($(end)+(mid)-(label)$);
   \node[] at ($(end)!7pt!(far)$) {$\scriptscriptstyle\cong$} ;
}

\newcommand{\sq}[5]{{#1}\colon({#4} \; ^{{#2}}_{\substack{{#3}}} \; {#5})}
\newcommand{\push}[3]{{#1}\,\square_{{#3}}\, {#2}}
\newcommand{\pushprod}[7]{{#1}\,\square\, {#4}\colon {#2}\otimes_{#7} {#6}\coprod_{{#2}\otimes_{#7} {#5}} {#3}\otimes_{#7} {#5}\to {#3}\otimes_{#7} {#6}}

\def\cellslide{0.5}
\def\celllength{.2cm}

\NewDocumentCommand{\cell}{ O{} O{n} O{\cellslide} O{\celllength} m m m }{
  \coordinate (mid) at ($({#5})!{#3}!({#6})$);
  \coordinate (start) at ($(mid)!{#4}!({#5})$);
  \coordinate (end) at ($(mid)!{#4}!({#6})$);
  \draw[#2] (start) to node
  [inner sep=6pt,outer sep=0,minimum size=0,#1]{{#7}} (end);
}


\DeclareMathOperator{\Loc}{Loc}